\section{Formal problem and the Hybrid-INoT core}
\label{sec:hybrid-core}

\subsection{Engineering task and observable success}
An engineering task is represented by $(x,\mathcal{C}_0,\mathcal{V})$: a specification $x$, a fixed supplied context $\mathcal{C}_0$, and an external verification procedure $\mathcal{V}$. The output consists of a candidate program $y$ and its verification record. Code generation and verification are different operations: the model proposes a solution; original executable tests determine the benchmark outcome in a controlled environment. A textual claim that the solution is correct cannot replace those tests. On repository tasks, $\mathcal{C}_0$ denotes the disclosed retrieved packet, not an assertion that the full repository was supplied.

For an execution protocol $a$, the engineering objective is to increase the probability of verified success while controlling expected token use and the estimated cost of generation. We therefore report quality and resources separately rather than optimize an unvalidated weighted score. Observed programs that fail tests are failures; unavailable generation or evaluation cannot be turned into a verified success or a fabricated test failure. Section~\ref{sec:heldout-methods} specifies the common eligibility controls and statistical endpoint.

\subsection{Internal role organization and external verification}
Hybrid-INoT uses a division of responsibility: planning, implementation and critique are prescribed inside a model response, while verification is performed outside the model. Here, \emph{hybrid} denotes this combination of internal instructions and external executable evidence. It does not require a learned deployment router. The original planner--worker--critic terminology corresponds to the executed labels \texttt{planner}, \texttt{implementer} and \texttt{reviewer}. Internal role segregation here denotes separation of instructed functions, not independently observed latent agents.

The single-call core studied here prescribes three operations: analyze requirements and edge cases, implement using the plan, then review and revise against the requirements. A single model call receives all three instructions and the complete supplied context and returns one final program. The experiment does not observe hidden plans or critic scores. The three operations are therefore instructed functions, not three measured internal subroutines. The precise prompts, including the final-output contract, appear in the appendix.

In the primary experiment, SR is this role-labelled single-call condition. SN is its role-label ablation with identical operation sentences. MR and MN distribute the corresponding operations across three fresh calls, and D removes the staged procedure altogether. All five receive the same subsequent native evaluation. The design tests whether the proposed internal organization contributes beyond a plain instruction and what changes when the same operations are distributed across call boundaries. Hybrid-INoT has no privileged access to verification.

\begin{figure}[htbp]\centering
\begin{tikzpicture}[
 x=1mm,y=1mm,
 box/.style={draw=black!80,align=center,font=\small,inner sep=2mm},
 flow/.style={-{Latex[length=2mm]},thick}]
\node[box,fill=gray!8,text width=100mm,minimum height=12mm] (input) at (0,0) {Task and complete supplied context};
\node[box,text width=43mm,minimum height=54mm] (single) at (-28,-40) {\textbf{SR: one call}\\[3mm]One prompt requests:\\[1mm]planning;\\implementation;\\requirements review\\[3mm]Internal steps\\are not measured};
\node[box,text width=43mm,minimum height=15mm] (m1) at (28,-20) {\textbf{MR: call 1}\\Prepare a plan};
\node[box,text width=43mm,minimum height=15mm] (m2) at (28,-40) {\textbf{MR: call 2}\\Write the program\\{\footnotesize All previous responses}};
\node[box,text width=43mm,minimum height=15mm] (m3) at (28,-60) {\textbf{MR: call 3}\\Review and revise\\{\footnotesize All previous responses}};
\node[box,text width=100mm,minimum height=14mm] (candidate) at (0,-84) {\textbf{Extract from the last response}\\One program; invalid format yields an empty candidate};
\node[box,fill=gray!8,text width=100mm,minimum height=18mm] (test) at (0,-110) {\textbf{External evaluation after generation}\\Original benchmark tests and common eligibility controls};
\node[box,text width=100mm,minimum height=20mm] (outcome) at (0,-139) {\textbf{Evaluation outcome}\\Pass, failure or unavailable evaluation\\Token use is recorded separately};
\draw[flow,dashed] (input.south) -- ++(0,-3) -| (single.north);
\draw[dashed,thick] (input.east) -- (61,0) -- (61,-60);
\foreach \call in {m1,m2,m3} {\draw[flow,dashed] (61,0 |- \call.east) -- (\call.east);}
\node[rotate=90,font=\footnotesize] at (65,-37) {Same task and context};
\draw[flow] (m1) -- (m2);
\draw[flow] (m2) -- (m3);
\draw[flow] (single.south) -- ++(0,-5) -| ([xshift=-28mm]candidate.north);
\draw[flow] (m3.south) -- ++(0,-5) -| ([xshift=28mm]candidate.north);
\draw[flow] (candidate) -- (test);
\draw[flow] (test) -- (outcome);
\end{tikzpicture}

\caption{Two generation configurations: SR requests three functions in one prompt; MR distributes them across three calls. Arrows between MR calls carry all complete previous responses; each call also receives the task and context. The figure shows the path of an observed response. Incomplete generation is handled separately in Algorithm~\ref{alg:hybrid-core}. All five conditions in Table~\ref{tab:design} share external evaluation, with no test feedback to the model.}
\label{fig:hybrid-core}
\end{figure}
\FloatBarrier

\subsection{State, execution and stopping}
The observable execution state is $\sigma_j=(x,\mathcal{C}_0,a,j,\mathcal{H}_j)$, where $a$ fixes the assigned condition, $j$ is the call index and $\mathcal{H}_j$ contains all complete preceding exposed responses. The model configuration is fixed within the comparison. D, SR and SN use $n_a=1$; MR and MN use $n_a=3$. In a multi-call condition, $\mathcal{H}_j$ is forwarded in full. The supplied context is unchanged: $\widehat{\mathcal{C}}=\mathcal{C}_0$. Condition assignment is fixed before generation, so it is not an outcome-dependent routing policy.

\begin{algorithm}[htbp]
\caption{Program generation and outcome evaluation}
\label{alg:hybrid-core}
\small
\textbf{Input:} task, complete supplied context and preassigned condition.\par
\textbf{Output:} observed program, evaluation outcome and known resource use.\par\medskip
\textbf{I. Generation}
\begin{enumerate}\setlength{\itemsep}{2pt}\setlength{\parsep}{0pt}
\item Determine the number of calls: one for D, SN or SR; three for MN or MR.
\item Execute the calls sequentially. Give each call the task, unchanged context and assigned instruction. Give calls two and three every complete previous response as well.
\item Retain each call's prompt, complete response, terminal status and token counters.
\item If a call is incomplete or resource accounting cannot be verified, retain the available evidence and stop that assignment as unavailable. Do not retry it automatically.
\item Extract one program from the last completed response. Retain an empty candidate if the format is invalid.
\end{enumerate}
\textbf{II. Evaluation after generation ends}
\begin{enumerate}\setcounter{enumi}{5}\setlength{\itemsep}{2pt}\setlength{\parsep}{0pt}
\item Evaluate every observed candidate with the original benchmark tests.
\item Match the evaluated program to its native report and common eligibility controls.
\item Record pass, failure or an unavailable outcome. Retain all known resource counters; do not return test results to the model.
\end{enumerate}
\end{algorithm}

Algorithm~\ref{alg:hybrid-core} describes the generation and evaluation stages of the experiment; the native harness is a subsequent evaluator, not a tool invoked inside the model call. Call deadlines and a finite $n_a$ bound generation attempts. This is a stopping rule, not a theorem of convergence or correctness. A rejected completion never triggers an automatic second attempt at the same cell. The disclosed continuation concerns only assignments for which no call had started.

This configuration fixes the original architecture's optional controls: context selection is the identity, the mode is assigned, the number of final candidates is one, and verification-driven reruns are disabled. These choices isolate the internal-role question from compression, a small checking model, and adaptive retries. They also define the scope of the evidence: the experiment evaluates this core, rather than every possible Hybrid-INoT deployment.
